
\section{Joy in Heaven}

(Luke 15:7)


“Folkebladet,” July 04, 1917

The question concerning God cannot readily be torn out of the human soul. Perhaps this comes entirely from the fact that it has come forth from God, created by him. Certainly it may be pressed back by other questions and interests, but it comes again—often strongly and irresistibly.

And what is it that we desire to know? What is there in the question concerning God that is so important for us helpless children of men to have made clear?

Yes, it is this: whether he cares for us, whether he has a loving heart toward us. And until we gain certainty upon this question, we shall never have peace. You who happen to read these words know that this is so. Nothing among all the earthly goods and gifts and earthly happiness you have attained can wholly satisfy you. Your soul longs for something more. And what is it? The answer is: God. Or perhaps more exactly: the certainty of being loved by God.

Am I, little and wretched as I am, evil and sinful as I am—am I loved by God? This is the fundamental question of life at the depth of the human heart, and upon this question I must have an answer. Who can give the answer? Other human beings? O yes, there are people who out of blessed experience can tell us that they know themselves to be loved by God, to have received grace, to be saved. It does us good to meet such people and to listen to their confession. But—it is not enough for me. For these people can be mistaken. Mistaken about God and about themselves. And what they tell me may rest upon a self-assumed delusion or something of that kind.

Is there, then, anyone who can give me an answer so certain and secure that I may dare to trust it, live upon it—die upon it?

Yes, there is. He who speaks in the Bible verse from which the heading of this little piece is taken—he can give an answer—an answer upon which I can rely in life and in death.

And he can give an answer because he knows God as no one else does. “The only-begotten Son, who is in the Father’s bosom, he has made him known.”

He gives his answer to the burning question of the human heart, whether God loves us, out of his own intimate fellowship and knowledge of his Father.

And his answer is this: “God so loved the world that he gave his Son!”

His answer is this: “I say to you: In the same way there will be joy in heaven over one sinner who repents!”

God is love, Jesus says. Thus we have received an answer to the question whether God cares for us. Jesus, his only-begotten Son, says: He loves us—loves, loves so that he gave his Son—gave him to the world, gave him to me!

Glorious truth! Blessed assurance: God is my Father.

I belong with him. And this Father does not forget his child—not even when the child forgets him—goes his own way and comes into misery.

This is the gospel. And the gospel is what the mission seeks to have proclaimed to as many as possible. The gospel is also what our Home Mission desires to have proclaimed to our people throughout the great cities and in the many new settlements where our people are building their homes.

You who happen to read this, will you do your part so that the gospel may be proclaimed to our people? Perhaps you yourself are so blessed as to know that for your sake there has been “joy in heaven”—when, repentant and bowed down, you stood before your God with this prayer: “Father, I have sinned!” and through the gospel received the assurance of full forgiveness of all sin. And is not this assurance the most glorious, the best of all?

O, so many of your brothers and sisters are going about in the world with a soul wounded by the arrow that came from the enemy. The gospel alone can heal the soul’s burning wounds—will you take part in sending it to our people?

At this time, throughout our whole country, collections are being made for “The Red Cross.” What is the purpose? To relieve bodily distress, to lessen pain and misery. Certainly a blessed work! But there is something in the world worse than bodily distress. You know this, my brother and sister, who in deep anguish of soul, in fear and trembling, have stood bowed before God in shame and sorrow over your sin. Everything else in the world seemed so insignificant to you then in comparison with the question of salvation, grace, mercy. And you know this also, that all the other joys and goods you possessed became small in comparison with what you felt, and still feel, when you were permitted to look into the Father’s heart, God’s great, warm, compassionate Father’s heart, and were enabled to believe that all is forgiven, all sin blotted out, remitted. Then you became glad in the midst of sorrow, rich in the midst of poverty, happy in the hardest adversity. And this happiness became yours through the gospel, the message of salvation from sin and death through Jesus Christ. Therefore I will repeat it: take part in sending the gospel out to our people in this country and in Canada! It is needed more than ever before that the message of God’s love for sinners sound forth clearly and in full voice. The time is troubled and difficult, and in many there is stirring a deep longing for something secure and steadfast to cling to:

“Where does my heart find its most blessed rest?
With Jesus, only with Jesus!
Where is found the firm-founded rock for faith?
With Jesus, yes, only with Jesus!”

My brother and sister, do what you can so that our people may be led to Jesus! Will you?

E. B.
